\chapter{安装与首次启动}\label{ch:install}
\index{Anzhuang yu Runtime@安装与 Runtime}

KSSOLV Toolbox 可以安装在自己的计算机上，也可以通过 Online 在浏览器中使用。两种方式提供相近的项目、建模和工作流操作。选择时，首先考虑课程或课题组已有的安排：如果教师已经提供 Online 入口，直接进入浏览器即可；如果希望在个人计算机上保存文件、离线准备结构和运行计算，可以安装独立应用程序。

本章从安装包的选择开始，介绍 Windows、macOS 和 Linux 的安装与首次启动。安装完成后，将用一个很小的练习检查界面和项目保存。Online 的登录、云端文件和会话操作见第~\ref{ch:online}~章。

\section{选择使用方式}

\subsection{独立应用程序与 MATLAB 工具箱}

独立应用程序使用 MATLAB Runtime 运行。Runtime 提供应用所需的运行环境，学生不必另外购买或安装完整的 MATLAB。它与 MATLAB 桌面程序是不同的软件；安装 Runtime 后，不会增加一个可供任意编程的 MATLAB 桌面。

如果计算机已经安装 MATLAB R2025a 或更新版本，也可以选择 MATLAB 工具箱。工具箱安装在 MATLAB 内，通过附加功能管理器或 MATLAB 中的启动入口打开。对日常图形操作而言，两者的主要工作流程相同；在调用工作区变量、自行编写脚本或使用其他 MATLAB 产品时，完整 MATLAB 环境更方便。

\begin{table}[htbp]
\centering
\caption{按使用环境选择安装方式}
\begin{tabularx}{\linewidth}{lX}
\toprule
使用方式 & 适合的情况 \\
\midrule
独立应用程序 & 希望直接从系统应用菜单打开工具箱，计算机没有完整 MATLAB。\\
MATLAB 工具箱 & 已有受支持的 MATLAB，并希望结合 MATLAB 工作区和脚本使用。\\
Online & 课程或课题组提供浏览器入口，希望使用统一配置的服务器环境。\\
\bottomrule
\end{tabularx}
\end{table}

\subsection{认识安装包中的 Runtime}

独立应用需要与该构建匹配的 MATLAB Runtime。安装包通常采用三种方式提供运行环境：联网包在需要时下载 Runtime，离线包随附 Runtime，而仅应用包使用计算机上已有的兼容 Runtime。应用版本号和 Runtime 的发行版号含义不同：前者标识 KSSOLV Toolbox，后者标识运行环境。

已经安装完整 MATLAB，并不意味着独立应用一定能直接找到所需的 Runtime。请让安装程序完成检测，按照窗口所示的版本和位置继续安装。也不要仅凭文件夹名称相似，选择另一发行版的 Runtime。

\begin{notebox}{选择适合当前计算机的文件}
下载时同时核对操作系统、处理器架构和安装包说明。Apple silicon 与 Intel Mac 使用不同架构的程序；Windows 和 Linux 的程序也不能互换。以发行页提供的文件和说明为准。
\end{notebox}

\section{安装前的准备}

从 KSSOLV Toolbox 项目发行页下载安装文件后，将文件放在容易找到的位置，例如“下载”文件夹。若课程提供了指定版本，优先使用教师给出的版本，以便与课堂中的菜单、示例和计算结果相对应。

安装前关闭正在运行的旧版 KSSOLV Toolbox。把需要保留的项目保存为 \file{.ks} 文件；升级通常保留项目和个人设置，但养成独立保存项目的习惯，也便于在另一台计算机继续工作。联网安装期间保持网络连接，离线安装时则检查随附文件是否齐全。

安装位置可选择有足够空间的本地磁盘。应用、Runtime 与个人项目可以放在不同目录：前两者由安装器管理，项目则建议集中放在一个自己熟悉的文件夹中。例如建立“计算化学课程”文件夹，再按练习或材料分别建立子文件夹。这样的安排比把所有项目留在下载目录更便于查找。

% SCREENSHOT: 发行文件列表与安装方式示意，突出操作系统、架构、联网/离线，而非构建信息。

\section{在 Windows 上安装}

\subsection{运行安装程序}

Windows 的独立分发包通常是带有版本号的 EXE 安装程序，文件名中包含 \file{KSSOLV_Toolbox} 和发行版本。部分离线包包含同名的 \file{.bin} 分卷文件。遇到这种形式时，把这些文件放在同一文件夹，再运行 EXE；不需要手动打开分卷。

\begin{steps}
\item 双击安装程序。若系统显示用户账户控制窗口，确认程序来源后允许继续。
\item 选择安装语言并阅读安装说明。简体中文安装界面适合与本手册对照使用。
\item 按页面要求确认工具箱和 Runtime 的许可，并继续到安装位置与 Runtime 设置。
\item 选择应用目录。可以保留默认位置，也可以选择另一块本地磁盘。
\item 检查 Runtime 的检测结果：已有兼容环境时按提示复用；没有时选择安装包提供的安装方式。
\item 点击安装，等待文件复制和最后的配置完成。显示完成页面后退出安装器。
\end{steps}

应用和 Runtime 的位置可以分别选择，例如：
\begin{itemize}[nosep]
\item 应用目录：\literalinput{E:/Apps/KSSOLV Toolbox}。
\item Runtime 目录：\literalinput{E:/MathWorks/MATLAB Runtime}。
\end{itemize}
当前安装流程要求 Runtime 目录使用 ASCII 字符，可使用英文、数字和空格；项目名称则可以按照自己的研究内容命名。

% SCREENSHOT: Windows 安装欢迎页；应用与 Runtime 目录页面，保留原生控件与中文说明。

\subsection{等待安装完成}

安装程序在最后阶段完成运行环境配置。此时进度可能比单纯复制文件慢，尤其是首次安装 Runtime 时。请等待安装器明确显示完成；不要因为文件已经出现在应用目录，就提前关闭安装程序。

如果页面提示配置没有完成，按提示重试，并检查磁盘空间、目录权限和 Runtime 来源。记录页面中的具体文字，通常比仅描述“安装失败”更有帮助。学校统一管理的计算机如要求管理员权限，可将提示交给负责该计算机的教师或管理员。

默认安装供所有用户使用，通常需要管理员授权。个人安装与已有 Runtime 的复用方式取决于分发包提供的选项。普通安装无需修改系统 PATH，也无需到 MATLAB 目录手动复制库文件。

\subsection{启动与更新}

完成安装后，从开始菜单或安装器创建的快捷方式启动 KSSOLV Toolbox。启动卡片会显示产品名称和当前启动状态；主窗口出现后，再等待项目面板和工具栏完成显示。

更新时运行新的安装包，按照向导替换旧版应用。已登记的旧版由安装器处理，共享 Runtime、个人配置和项目通常保留。应用安装目录不应当作为个人项目目录，后续升级时也更容易分清程序文件与自己的数据。

\begin{notebox}{离线包的检查点}
如果安装器提示找不到附加文件，先检查 EXE 和全部同名分卷是否在一起，文件名是否被浏览器自动增加了重复编号。确认下载完整后，再重新运行安装程序。
\end{notebox}

\section{在 macOS 上安装}

\subsection{选择与 Mac 匹配的分发包}

在苹果菜单的“关于本机”中可以查看处理器信息。当前 Apple silicon 安装包面向采用 Apple 芯片的 Mac。下载时核对文件说明，不要仅根据文件名中出现“macOS”便判断架构适用。

macOS 使用原生安装器完成 PKG 安装。应用通常放在 \file{/Applications}，安装器会检测所需的 Runtime，并根据安装包类型复用、下载或安装运行环境。整个过程在系统安装向导中完成。

安装向导的介绍页列出应用版本、处理器架构和 MATLAB Runtime 要求。开始安装前，核对这些信息是否与当前 Mac 匹配，并确认安装包是仅包含应用，还是同时提供运行环境；具体要求以所用安装包的介绍页为准。

\begin{steps}
\item 打开下载的分发文件，运行其中的 PKG 安装包。
\item 阅读介绍页，确认应用版本、处理器架构和 Runtime 发行版。
\item 按向导阅读并确认许可，然后进入安装类型与位置设置。
\item 选中 KSSOLV Toolbox，检查应用的安装位置；通常保留“应用程序”目录即可。
\item 检查 Runtime 设置，按下一小节说明选择已有环境或由安装器准备环境。
\item 开始安装，并在系统要求时完成管理员认证。
\item 等待安装程序报告成功，再从“应用程序”打开 KSSOLV Toolbox。
\end{steps}

许可页如图~\ref{fig:install-macos-license}~所示。可直接阅读条款，也可使用窗口下方的 \ui{打印…} 或 \ui{保存…} 留存文本。阅读后点击 \ui{继续}，按随后出现的对话框确认是否同意；同意后才能进入后续安装步骤。向导最后显示成功页面后，安装才告完成。

\screenshot{publication-assets/crop-b9b8c6192e335985ecb0.png}{macOS 安装向导的许可页。图中正在阅读软件许可，尚未完成许可确认，也尚未执行安装。}{fig:install-macos-license}

% SCREENSHOT: macOS 安装类型/Runtime位置页与安装成功页，待完成相应实机步骤后补入；不要将许可页当作安装完成证据。

\subsection{处理已有的 Runtime}

如果使用仅应用安装包，计算机上必须已经具备兼容的 Runtime。自动检测不到时，在安装类型页选中“已有 MATLAB Runtime”，通过“位置”进入“其他…”，选择安装器所要求的发行版文件夹。应选择 Runtime 的根目录，而不是它内部某个库目录。

联网安装包会在缺少 Runtime 时从 MathWorks 下载；离线包则使用随包提供的内容。已有环境的发行版、架构和所需组件符合要求时可以复用。安装器完成 Runtime 准备后，再替换旧版工具箱。

安装最后阶段会在后台完成配置，不一定出现 Toolbox 主窗口。待系统安装器显示成功后，再主动启动应用。第一次打开时出现启动卡片属于正常过程；后续每次启动仍需要加载运行环境，只是无需重复进行安装。

\subsection{保留项目并升级}

升级前先保存并关闭应用，再运行新安装包。共享 Runtime、用户配置和项目不需要因为更换 Toolbox 而逐一移走。建议把 \file{.ks} 文件放在自己的“文稿”或课题目录，而不是应用包内部。

如果系统询问是否允许打开从互联网下载的应用，按照 macOS 显示的来源信息和单位的软件管理要求操作。遇到由学校管理策略限制的情况，应将系统的实际提示交给管理员；无需为启动工具箱关闭整个系统的安全设置。

\section{在 Linux 上安装}

Linux 独立应用应使用对应平台的安装文件，并在图形桌面会话中运行。下载完成后，若文件已经可以直接执行，可以从文件管理器启动；如果没有执行权限，可在文件所在目录打开终端。

\begin{steps}
\item 在终端进入安装文件所在目录。
\item 为安装文件增加执行权限：\texttt{chmod +x KSSOLV\_Toolbox.install}。
\item 运行 \texttt{./KSSOLV\_Toolbox.install}，进入图形安装向导。
\item 按向导选择应用和 Runtime 位置，等待安装完成。
\item 按安装完成页提供的入口或启动命令启动应用。
\end{steps}

这里的文件名用于说明常见形式，具体名称以下载内容为准。安装程序最后显示的 Runtime 路径与启动命令应保存下来，尤其是在没有自动创建应用菜单项的桌面环境中。

若在远程终端中只看到命令行而没有图形桌面，应先确认是否准备使用 Linux 桌面应用，还是由单位提供的 Online。对普通浏览器用户来说，不需要在自己的计算机运行 Linux 安装命令；打开 Online 入口即可。

% SCREENSHOT: Linux 安装目录/Runtime 选择窗口和安装完成页，采用真实目标发行版。

\section{在 MATLAB 中安装工具箱}

已有完整 MATLAB 的读者可以使用 \file{KSSOLV_Toolbox.mltbx}。开始前，在 MATLAB 中确认版本达到 R2025a 或更新版本，然后完成以下操作。

\begin{steps}
\item 从项目发行页下载 \file{KSSOLV_Toolbox.mltbx}。
\item 在 MATLAB 中打开该文件，按附加功能安装窗口完成安装。
\item 从 MATLAB 或附加功能管理器启动 KSSOLV Toolbox。
\item 确认主窗口出现，再进行项目练习。
\end{steps}

如果计算机上已有其他版本，使用 MATLAB 的附加功能管理器查看和管理已安装的工具箱。课堂中最好统一 MATLAB 和 Toolbox 的主要版本；这样在讨论菜单和示例时，可以把注意力集中在计算方法上。

本手册以图形界面为主。后面的结构导入、建模和工作流步骤，既适用于独立应用，也适用于 MATLAB 中的工具箱；只有明确提到 MATLAB 工作区或额外 MATLAB 产品时，才需要相应环境。

\section{第一次启动后的练习}

打开独立应用程序时，会先出现启动卡片。图~\ref{fig:desktop-startup-050}~中的卡片显示产品版本与当前加载状态，此时应用正在准备工作区。首次启动或更新后，加载运行环境可能需要较长时间；保留启动窗口，等待主窗口出现。

\begin{figure}[!htb]
\centering
\includegraphics[width=122mm]{publication-assets/crop-6bf7d3c41caf0892911f.png}
\caption{macOS 中 KSSOLV Toolbox 0.5.0 独立应用的实际启动卡片。图中显示“正在准备变量数据…”，属于进入工作区前的加载状态。}
\label{fig:desktop-startup-050}
\end{figure}

启动卡片收起后，主窗口提供项目浏览器、功能区和中央文档区。等待这些区域完成显示，再开始下面的版本检查与项目练习。启动过程中的状态文字描述应用的准备工作；电子结构计算则由进入工作区后建立的工作流执行。

\subsection{确认界面和版本}

主窗口由上方工具栏、侧边浏览面板和中央文档区组成。第一次进入时不必立即理解所有按钮；先找到 \ui{主页}、\ui{设置} 和 \ui{资源}，再观察项目浏览器的位置。

选择 \ui{主页 → 帮助 → 关于 KSSOLV}，核对产品版本。与课程截图存在小幅差别时，先比较版本和界面语言。语言可在 \ui{主页 → 设置 → 通用} 中选择，修改后下次启动生效。

\subsection{创建一个可以重新打开的项目}

这项练习只检查文件与界面，不需要运行计算。准备随书的 \file{examples/silicon.cif}，或课程提供的另一个小型 CIF、XYZ 结构文件。使用随书硅结构，可以直接与下一章的文件练习和工作区介绍衔接。

\begin{steps}
\item 选择 \ui{主页 → 结构 → 从文件导入结构…}，选中课程结构文件。
\item 确认结构名称出现在项目浏览器中，中央三维视图显示原子。
\item 用鼠标拖动旋转视图，再用滚轮调整大小，确认结构能够完整显示。
\item 点击 \ui{保存}，在自己的课程目录输入 \file{first-project.ks}。
\item 确认保存完成后关闭当前项目，再通过 \ui{打开项目…} 打开刚保存的文件。
\item 检查项目浏览器和结构文档，确认此前导入的结构仍在。
\end{steps}

% SCREENSHOT: 首次导入后的主窗口与保存文件窗口，强调项目树、三维结构和 .ks 扩展名。

完成这项练习，说明安装、图形显示和项目基本读写已经可以使用。它尚未检查具体材料的计算参数；第一次电子结构计算见第~\ref{ch:first}~章。

\section{遇到问题时先检查什么}

如果程序未能启动，先记录 Toolbox 版本、操作系统，以及安装或启动窗口中的完整提示。再检查所用安装包是否匹配处理器架构，Runtime 是否为安装器要求的版本。将提示原文与安装位置一并提供，比反复重新安装更有助于判断原因。

如果主窗口出现但三维结构没有显示，先用本章的小结构重试，并观察是否有明确的图形提示。Online 还需使用能运行 WebGL 的浏览器。若文件导入后没有新增项目项，则核对结构文件格式和文件内容，不必立即判断为安装失败。

如果保存时找不到文件，留意自己使用的是桌面应用还是 Online：桌面选择本机目录，正式 Online 通常保存在云端“文件”面板。下一章将专门介绍这两种保存位置的区别，并完成一次从上传到下载的完整练习。

\subsection{为后续练习整理文件}

为课程或课题建立固定目录，分别保存原始结构和工作项目，保留原始文件的来源与名称。在结构准备、参数检查和结果整理完成时另存项目副本，便于回到已有阶段。跨计算机传递时优先使用完整的 \file{.ks} 项目；单独的结构文件通常不包含工作流和计算结果。

